\chapter{2005年北京卷（春文）}
\section{单选题}
\begin{problem}
\(\i-2\) 的共轭复数是（\(\quad\)）

\choices
    {\(2+\i\)}
    {\(2-\i\)}
    {\(-2+\i\)}
    {\(-2-\i\)}
\end{problem}

\begin{answer}
D
\end{answer}

\begin{solution}
令 \(z=\i-2=-2+\i\). 共轭复数只改变虚部的符号，因此 \(\overline z=-2-\i\)，故选 D.
\end{solution}

\begin{problem}
函数 \(y=|\log_2x|\) 的图象是（\(\quad\)）

\choices
    {\choicetexfigure[width=0.15\paperwidth]{img_repaint/2005/beijing_liberal_spring/q2_optA.tex}}
    {\choicetexfigure[width=0.15\paperwidth]{img_repaint/2005/beijing_liberal_spring/q2_optB.tex}}
    {\choicetexfigure[width=0.15\paperwidth]{img_repaint/2005/beijing_liberal_spring/q2_optC.tex}}
    {\choicetexfigure[width=0.15\paperwidth]{img_repaint/2005/beijing_liberal_spring/q2_optD.tex}}
\end{problem}

\begin{answer}
A
\end{answer}

\begin{solution}
由对数真数的条件，函数的定义域为 \(x>0\). 在此定义域内，按 \(\log_2x\) 的正负分段，得
\[
y=|\log_2x|=\begin{cases}
-\log_2x,&0<x<1,\\
\log_2x,&x\geq 1.
\end{cases}
\]
当 \(0<x<1\) 时，\(\log_2x<0\)，所以图象是 \(y=\log_2x\) 在 \(x\) 轴下方部分关于 \(x\) 轴的对称图象；当 \(x\geq1\) 时，图象就是 \(y=\log_2x\). 因而图象经过 \((1,0)\)，在该点取得最小值，左支单调递减，右支单调递增，故选 A.
\end{solution}

\begin{problem}
下列命题中，正确的是（\(\quad\)）

\choices
    {经过不同的三点有且只有一个平面}
    {分别在两个平面内的两条直线一定是异面直线}
    {垂直于同一个平面的两条直线是平行直线}
    {垂直于同一个平面的两个平面平行}
\end{problem}

\begin{answer}
C
\end{answer}

\begin{solution}
命题（A）不正确，因为三点可以共线；此时经过这条直线的平面有无数个，不是唯一一个平面.

命题（B）不正确. 例如取两个相交平面，分别在两个平面内取经过它们交线同一点的两条直线，这两条直线可以相交，所以不一定是异面直线.

命题（C）正确. 若两条直线都垂直于平面 \(\alpha\)，则它们都平行于平面 \(\alpha\) 的法线，因而两条直线互相平行.

命题（D）不正确. 例如平面 \(xOy\) 与平面 \(yOz\)、平面 \(xOz\) 都垂直，但后两个平面相交于 \(z\) 轴，所以它们不一定平行. 因此正确命题为（C），故选 C.
\end{solution}

\begin{problem}
如果函数 \(f(x)=\sin(\pi x+\theta)\)（\(0<\theta<2\pi\)）的最小正周期是 \(T\)，且当 \(x=2\) 时取得最大值，那么（\(\quad\)）

\choices
    {\(T=2\)，\(\theta=\dfrac{\pi}{2}\)}
    {\(T=1\)，\(\theta=\pi\)}
    {\(T=2\)，\(\theta=\pi\)}
    {\(T=1\)，\(\theta=\dfrac{\pi}{2}\)}
\end{problem}

\begin{answer}
A
\end{answer}

\begin{solution}
对于函数 \(\sin(\omega x+\theta)\)，最小正周期为 \(\dfrac{2\pi}{|\omega|}\). 这里 \(\omega=\pi\)，所以 \(T=2\). 当 \(x=2\) 时取得最大值，必须有
\[
\sin(2\pi+\theta)=\sin\theta=1
\]
因此 \(\theta=\dfrac{\pi}{2}+2k\pi\). 结合 \(0<\theta<2\pi\)，只能取 \(\theta=\dfrac{\pi}{2}\)，故选 A.
\end{solution}

\begin{problem}
设 \(abc\ne0\)，“\(ac>0\)”是“曲线 \(ax^2+by^2=c\) 为椭圆”的（\(\quad\)）

\choices
    {充分非必要条件}
    {必要非充分条件}
    {充分必要条件}
    {既非充分又非必要条件}
\end{problem}

\begin{answer}
B
\end{answer}

\begin{solution}
若曲线 \(ax^2+by^2=c\) 是椭圆，由于 \(a\)，\(b\)，\(c\ne0\)，将等式除以 \(c\) 后，\(x^2\) 和 \(y^2\) 的系数必须同为正，故 \(a\)，\(b\)，\(c\) 必须同号，从而必有 \(ac>0\). 所以 \(ac>0\) 是必要条件.

但它不是充分条件. 例如取 \(a=c=1\)，\(b=-1\)，则 \(ac=1>0\)，而曲线变为 \(x^2-y^2=1\)，这是双曲线，不是椭圆. 因此 \(ac>0\) 是必要非充分条件，故选 B.
\end{solution}

\begin{problem}
直线 \(x+\sqrt{3}y-2=0\) 被圆 \((x-1)^2+y^2=1\) 所截得的线段的长为（\(\quad\)）

\choices
    {\(1\)}
    {\(\sqrt{2}\)}
    {\(\sqrt{3}\)}
    {\(2\)}
\end{problem}

\begin{answer}
C
\end{answer}

\begin{solution}
圆心为 \(C(1,0)\)，半径为 \(r=1\). 圆心到直线的距离为
\[
d=\frac{|1+\sqrt{3}\cdot0-2|}{\sqrt{1+(\sqrt{3})^2}}=\frac12
\]
设弦的中点为圆心到直线的垂足，则圆心、弦中点和弦端点组成直角三角形，故截得弦长为
\[
2\sqrt{1^2-d^2}=2\sqrt{1-\frac14}=\sqrt{3}
\]
故选 C.
\end{solution}

\begin{problem}
在 \(\triangle ABC\) 中，已知 \(2\sin A\cos B=\sin C\)，那么 \(\triangle ABC\) 一定是（\(\quad\)）

\choices
    {直角三角形}
    {等腰三角形}
    {等腰直角三角形}
    {正三角形}
\end{problem}

\begin{answer}
B
\end{answer}

\begin{solution}
由三角形内角和，\(C=\pi-A-B\)，于是
\[
\sin C=\sin(A+B)=\sin A\cos B+\cos A\sin B
\]
代入题设 \(2\sin A\cos B=\sin C\)，得 \(\sin A\cos B=\cos A\sin B\)，即
\[
\sin A\cos B-\cos A\sin B=\sin(A-B)=0
\]
因为 \(0<A<\pi\)，\(0<B<\pi\)，所以 \(-\pi<A-B<\pi\)，在此范围内 \(\sin(A-B)=0\) 只能推出 \(A-B=0\)，即 \(A=B\). 因而 \(A\)，\(B\) 所对的边 \(BC\)，\(CA\) 相等，三角形为等腰三角形，故选 B.
\end{solution}

\begin{problem}
若不等式 \((-1)^na<2+\dfrac{(-1)^{n+1}}{n}\) 对于任意正整数 \(n\) 恒成立，则实数 \(a\) 的取值范围是（\(\quad\)）

\choices
    {\(\left[-2,\dfrac{3}{2}\right)\)}
    {\(\left(-2,\dfrac{3}{2}\right)\)}
    {\(\left[-3,\dfrac{3}{2}\right)\)}
    {\(\left(-3,\dfrac{3}{2}\right)\)}
\end{problem}

\begin{answer}
A
\end{answer}

\begin{solution}
当 \(n=2m\) 为偶数时，\((-1)^n=1\)，\((-1)^{n+1}=-1\)，所以不等式化为
\[
a<2-\frac{1}{2m}
\]
对所有正整数 \(m\)，右端的最小值在 \(m=1\) 时为 \(\dfrac32\)，故偶数情形等价于 \(a<\dfrac32\).

当 \(n=2m-1\) 为奇数时，\((-1)^n=-1\)，\((-1)^{n+1}=1\)，所以
\[
-a<2+\frac{1}{2m-1}
\]
即 \(a>-2-\dfrac{1}{2m-1}\). 若 \(a\geq-2\)，这个不等式对所有 \(m\) 都成立. 若 \(a<-2\)，令 \(\varepsilon=-a-2>0\)，取正整数 \(m\) 使 \(2m-1>\dfrac1\varepsilon\)，则 \(\dfrac1{2m-1}<\varepsilon\)，从而
\[
-2-\frac1{2m-1}>-2-\varepsilon=a
\]
这与所需的 \(a>-2-\dfrac1{2m-1}\) 矛盾. 因此奇数情形等价于 \(a\geq-2\).

综上，\(a\) 的取值范围为 \(\left[-2,\dfrac32\right)\)，故选 A.
\end{solution}

\section{填空题}
\begin{problem}
\(\displaystyle\lim_{n\to\infty}\dfrac{n^2}{2n^2-3}=\)\fillinblank{}.
\end{problem}

\begin{answer}
\(\dfrac12\)
\end{answer}

\begin{solution}
因为 \(n\ne0\)，分子、分母同除以 \(n^2\)，有
\[
\frac{n^2}{2n^2-3}=\frac{1}{2-\frac{3}{n^2}}
\]
当 \(n\to\infty\) 时，\(\dfrac{3}{n^2}\to0\)，且分母的极限为 \(2\ne0\)，所以
\[
\displaystyle\lim_{n\to\infty}\frac{n^2}{2n^2-3}=\frac{1}{2-0}=\frac12
\]
\end{solution}

\begin{problem}
椭圆 \(\dfrac{x^2}{25}+\dfrac{y^2}{9}=1\) 的离心率是\fillinblank{}，准线方程是\fillinblank{}.
\end{problem}

\begin{answer}
\(\dfrac45\)，\(x=\pm\dfrac{25}{4}\)
\end{answer}

\begin{solution}
将方程与标准形式比较，长半轴为 \(a=5\)，短半轴为 \(b=3\)，焦半距满足
\[
c^2=a^2-b^2=25-9=16
\]
所以 \(c=4\). 因而离心率为 \(e=\dfrac ca=\dfrac45\). 椭圆焦点在 \(x\) 轴上，准线方程为 \(x=\pm\dfrac ae\)，故
\[
x=\pm\frac{5}{4/5}=\pm\frac{25}{4}
\]
\end{solution}

\begin{problem}
已知 \(\sin\dfrac{\theta}{2}+\cos\dfrac{\theta}{2}=\dfrac{2\sqrt{3}}{3}\)，那么 \(\sin\theta\) 的值为\fillinblank{}，\(\cos2\theta\) 的值为\fillinblank{}.
\end{problem}

\begin{answer}
\(\dfrac13\)，\(\dfrac79\)
\end{answer}

\begin{solution}
将已知等式两边平方，得
\[
\sin^2\frac\theta2+\cos^2\frac\theta2+2\sin\frac\theta2\cos\frac\theta2=\frac43
\]
利用 \(\sin^2\dfrac\theta2+\cos^2\dfrac\theta2=1\) 以及 \(2\sin\dfrac\theta2\cos\dfrac\theta2=\sin\theta\)，可得 \(1+\sin\theta=\dfrac43\)，所以 \(\sin\theta=\dfrac13\). 再用二倍角公式，
\[
\cos2\theta=1-2\sin^2\theta=1-2\left(\frac13\right)^2=\frac79
\]
\end{solution}

\begin{problem}
如图，正方体 \(ABCD-A_1B_1C_1D_1\) 的棱长为 \(a\)，将该正方体沿对角面 \(BB_1D_1D\) 切成两块，再将这两块拼接成一个不是正方体的四棱柱，那么所得四棱柱的全面积为\fillinblank{}.

\texfigure[max width=0.42\linewidth]{img_repaint/2005/beijing_liberal_spring/q12_fig1.tex}
\end{problem}

\begin{answer}
\((4+2\sqrt{2})a^2\)
\end{answer}

\begin{solution}
切面 \(BB_1D_1D\) 的底边为正方形的对角线，切开后得到两个全等的直三棱柱. 每个三棱柱的三角形底面为直角等腰三角形，直角边均为 \(a\)，斜边为 \(a\sqrt2\). 因此两个三角形底面的面积之和为
\[
2\cdot\frac12a^2=a^2
\]
三个侧面的面积之和为
\[
a\cdot a+a\cdot a+a\cdot a\sqrt2=(2+\sqrt2)a^2
\]
所以一个三棱柱的表面积为 \((3+\sqrt2)a^2\).

若沿切面 \(BB_1D_1D\) 对应的矩形侧面相接，该侧面的面积为 \(a^2\sqrt2\)，两个三角形底面恰好拼成正方形 \(ABCD\)，所得棱柱就是原正方体；由于题设要求所得四棱柱不是正方体，不能沿此矩形侧面相接，只能沿面积为 \(a^2\) 的正方形侧面相接. 相接的两个面变为内部面，不再计入外表面积. 所得四棱柱的全面积为
\[
2(3+\sqrt2)a^2-2a^2=(4+2\sqrt2)a^2
\]
\end{solution}

\begin{problem}
从 \(0\)，\(1\)，\(2\)，\(3\) 这四个数中选三个不同的数作为函数 \(f(x)=ax^2+bx+c\) 的系数，可组成不同的一次函数有\fillinblank{}个，不同的二次函数有\fillinblank{}个.（用数字作答）
\end{problem}

\begin{answer}
\(6\)，\(18\)
\end{answer}

\begin{solution}
一次函数要求二次项系数 \(a=0\). 此时 \(b\)，\(c\) 必须从 \(1\)，\(2\)，\(3\) 中选取两个不同的数，并且分别放在 \(b\)，\(c\) 的位置上，因此有 \(3\times2=6\) 个.

二次函数要求 \(a\ne0\). 先从 \(1\)，\(2\)，\(3\) 中选取 \(a\)，有 \(3\) 种；确定 \(a\) 后，还剩下 \(0\) 以及另外两个非零数共三个数，从中有序选取两个分别作为 \(b\)，\(c\)，有 \(3\times2\) 种. 因此不同的二次函数有 \(3\times3\times2=18\) 个.
\end{solution}

\begin{problem}
若关于 \(x\) 的不等式 \(x^2-ax-a>0\) 的解集为 \((-\infty,+\infty)\)，则实数 \(a\) 的取值范围是\fillinblank{}.
\end{problem}

\begin{answer}
\((-4,0)\)
\end{answer}

\begin{solution}
令 \(h(x)=x^2-ax-a\). 这是开口向上的二次函数. 要使不等式的解集为全体实数，必须使 \(h(x)\) 对任意实数都严格为正；若判别式等于零，函数会在重根处取零，因此必须且只需
\[
\Delta=(-a)^2-4\cdot1\cdot(-a)=a^2+4a<0
\]
解不等式得 \(a(a+4)<0\)，所以 \(-4<a<0\). 因此实数 \(a\) 的取值范围为 \((-4,0)\).
\end{solution}

\section{解答题}
\begin{problem}
设函数 \(f(x)=\lg(2x-3)\) 的定义域为集合 \(M\)，函数 \(g(x)=\sqrt{(x-3)(x-1)}\) 的定义域为集合 \(N\). 求：

\begin{enumerate}
\item 集合 \(M\)，\(N\)；
\item 集合 \(M\cap N\)，\(M\cup N\).
\end{enumerate}
\end{problem}

\begin{solution}
\begin{enumerate}
\item 对数的真数必须为正，因此 \(2x-3>0\)，解得
\[
M=\left(\frac32,+\infty\right)
\]
根式内的式子必须非负，因此
\[
(x-3)(x-1)\geq0
\]
临界点为 \(1,3\)，作符号判断可得
\[
N=(-\infty,1]\cup[3,+\infty)
\]
\item 交集要求同时属于 \(M\) 和 \(N\). 由于 \(M\) 中的数大于 \(\dfrac32\)，只有 \(N\) 的右半部分与之相交，故
\[
M\cap N=[3,+\infty)
\]
并集把两个集合的所有部分合并，故
\[
M\cup N=(-\infty,1]\cup\left(\frac32,+\infty\right)
\]
\end{enumerate}
\end{solution}

\begin{problem}
如图，正三棱锥 \(S-ABC\) 中，底面边长是 \(3\)，棱锥的侧面积等于底面积的 \(2\) 倍，\(M\) 是 \(BC\) 的中点. 求：

\begin{enumerate}
\item \(\dfrac{AM}{SM}\) 的值；
\item 二面角 \(S-BC-A\) 的大小；
\item 正三棱锥 \(S-ABC\) 的体积.
\end{enumerate}

\texfigure[max width=0.48\linewidth]{img_repaint/2005/beijing_liberal_spring/q16_fig1.tex}
\end{problem}

\begin{solution}
设顶点 \(S\) 在底面上的射影为正三角形中心 \(O\). 底面是边长为 \(3\) 的正三角形，面积为 \(\dfrac{9\sqrt3}{4}\). 侧面积是底面积的两倍，且三个侧面全等，所以每个侧面的面积为
\[
\frac13\cdot2\cdot\frac{9\sqrt3}{4}=\frac{3\sqrt3}{2}
\]

\begin{enumerate}
\item 在正三角形 \(ABC\) 中，中线 \(AM\) 也是高，所以 \(AM=\dfrac{3\sqrt3}{2}\). 侧面 \(SBC\) 是等腰三角形，且 \(M\) 是 \(BC\) 的中点，因此 \(SM\perp BC\)，\(SM\) 是该侧面的高. 由侧面面积得
\[
\frac12\cdot BC\cdot SM=\frac12\cdot3\cdot SM=\frac{3\sqrt3}{2}
\]
从而 \(SM=\sqrt3\). 因此
\[
\frac{AM}{SM}=\frac{3\sqrt3/2}{\sqrt3}=\frac32
\]
\item 正三角形的重心、外心、内心重合，所以 \(A\)，\(O\)，\(M\) 共线，且重心把中线按顶点到中点的方向分成 \(2:1\). 因此
\(OM=\frac13AM=\frac{\sqrt3}{2}\)，\(AO=\frac23AM=\sqrt3\).
由于 \(SO\perp\) 底面，\(SO\perp OM\)，在直角三角形 \(SOM\) 中
\[
SO=\sqrt{SM^2-OM^2}=\sqrt{3-\frac34}=\frac32
\]
又因为 \(SO\perp AO\)，在直角三角形 \(SOA\) 中
\[
AS^2=SO^2+AO^2=\frac94+3=\frac{21}{4}
\]

二面角 \(S-BC-A\) 的平面角是直线 \(SM\) 与 \(AM\) 的夹角. 在三角形 \(SAM\) 中，设该角为 \(\varphi\)，则
\[
\cos\varphi=\frac{SM^2+AM^2-AS^2}{2\cdot SM\cdot AM}=\frac{3+\frac{27}{4}-\frac{21}{4}}{2\cdot\sqrt3\cdot\frac{3\sqrt3}{2}}=\frac12
\]
因为 \(0^\circ<\varphi<180^\circ\)，所以 \(\varphi=60^\circ\). 这里 \(SM\) 和 \(AM\) 都垂直于棱 \(BC\)，故它们的夹角就是二面角 \(S-BC-A\) 的平面角.
\item 正三棱锥的高为 \(SO=\dfrac32\)，所以体积为
\[
V=\frac13\cdot\frac{9\sqrt3}{4}\cdot SO=\frac13\cdot\frac{9\sqrt3}{4}\cdot\frac32=\frac{9\sqrt3}{8}
\]
\end{enumerate}
\end{solution}

\begin{problem}
已知 \(\{a_n\}\) 是等比数列，\(a_1=2\)，\(a_4=54\)；\(\{b_n\}\) 是等差数列，\(b_1=2\)，\(b_1+b_2+b_3+b_4=a_1+a_2+a_3\).

\begin{enumerate}
\item 求数列 \(\{a_n\}\) 的通项公式及前 \(n\) 项和 \(S_n\) 的公式；
\item 求数列 \(\{b_n\}\) 的通项公式；
\item 设 \(U_n=b_1+b_4+b_7+\cdots+b_{3n-2}\)，其中 \(n=1\)，\(2\)，\(\cdots\)，求 \(U_{10}\) 的值.
\end{enumerate}
\end{problem}

\begin{solution}
\begin{enumerate}
\item 设等比数列 \(\{a_n\}\) 的公比为 \(q\). 由 \(a_1=2\)、\(a_4=54\)，有
\(a_4=a_1q^3\)，\(\Longrightarrow\)，\(2q^3=54\)，\(\Longrightarrow\)，\(q^3=27\).
由于 \(q\) 为实数，得 \(q=3\). 所以
\[
a_n=a_1q^{n-1}=2\cdot3^{n-1}
\]
前 \(n\) 项和为
\[
S_n=a_1\frac{q^n-1}{q-1}=2\cdot\frac{3^n-1}{2}=3^n-1
\]
\item 由上问，\(a_1+a_2+a_3=2+6+18=26\). 设等差数列 \(\{b_n\}\) 的公差为 \(d\)，则 \(b_4=2+3d\)，四项和为
\[
b_1+b_2+b_3+b_4=\frac{4(b_1+b_4)}2=\frac{4(2+2+3d)}2=8+6d
\]
由题设 \(8+6d=26\)，解得 \(d=3\). 因此
\[
b_n=b_1+(n-1)d=2+3(n-1)=3n-1
\]
\item 按题目给出的下标，\(b_1\)，\(b_4\)，\(b_7\)，\(\ldots\)，\(b_{3n-2}\) 构成首项为 \(b_1=2\)、公差为 \(3d=9\) 的等差数列. 当 \(n=10\) 时，末项下标为 \(3\times10-2=28\)，故
\[
b_{28}=3\cdot28-1=83
\]
所以
\[
U_{10}=\frac{10(2+83)}{2}=425
\]
\end{enumerate}
\end{solution}

\begin{problem}
如图，\(O\) 为坐标原点，过点 \(P(2,0)\) 且斜率为 \(k\) 的直线 \(l\) 交抛物线 \(y^2=2x\) 于 \(M(x_1,y_1)\)，\(N(x_2,y_2)\) 两点.

\begin{enumerate}
\item 写出直线 \(l\) 的截距式方程；
\item 求 \(x_1x_2\) 与 \(y_1y_2\) 的值；
\item 求证：\(OM\perp ON\).
\end{enumerate}

\texfigure[max width=0.42\linewidth]{img_repaint/2005/beijing_liberal_spring/q18_fig1.tex}
\end{problem}

\begin{solution}
直线 \(l\) 过 \(P(2,0)\)，且与抛物线交于两个点. 若 \(k=0\)，则 \(l\) 为 \(x\) 轴，与 \(y^2=2x\) 只有一个交点 \(O\)，故 \(k\ne0\). 由点斜式，直线方程为 \(y=k(x-2)\).

\begin{enumerate}
\item 令 \(y=0\)，得直线在 \(x\) 轴上的截距为 \(2\)；令 \(x=0\)，得直线在 \(y\) 轴上的截距为 \(-2k\). 因此截距式方程为
\[
\frac{x}{2}+\frac{y}{-2k}=1
\]
即 \(\dfrac{x}{2}-\dfrac{y}{2k}=1\).
\item 由直线方程得 \(x=2+\dfrac{y}{k}\). 代入抛物线方程，得
\[
y^2=2\left(2+\frac{y}{k}\right)
\]
即
\[
y^2-\frac{2}{k}y-4=0
\]
由于 \(M\)，\(N\) 是直线与抛物线的两个交点，\(y_1\)，\(y_2\) 是这个关于 \(y\) 的二次方程的两个根. 由韦达定理，
\[
y_1y_2=-4
\]
又由抛物线方程 \(x_i=\dfrac{y_i^2}{2}\)，所以
\[
x_1x_2=\frac{y_1^2}{2}\cdot\frac{y_2^2}{2}=\frac{(y_1y_2)^2}{4}=4
\]
\item 由 \(y_1y_2=-4\ne0\)，可知 \(y_1\)，\(y_2\ne0\)，从而 \(x_1\)，\(x_2>0\)，直线 \(OM\)，\(ON\) 的斜率存在. 利用 \(x_i=\dfrac{y_i^2}{2}\)，有
\(k_{OM}=\frac{y_1}{x_1}=\frac{2}{y_1}\)，\(k_{ON}=\frac{y_2}{x_2}=\frac{2}{y_2}\).
因此
\[
k_{OM}k_{ON}=\frac{4}{y_1y_2}=-1
\]
两条非竖直直线斜率之积为 \(-1\)，所以 \(OM\perp ON\).
\end{enumerate}
\end{solution}

\begin{problem}
经过长期观测得到：在交通繁忙的时段内，某公路段汽车的车流量 \(y\)（\(\text{千辆/小时}\)）与汽车的平均速度 \(v\)（\(\text{千米/小时}\)）之间的函数关系为：
\(y=\dfrac{920v}{v^2+3v+1600}\)（\(v>0\)）.

\begin{enumerate}
\item 在该时段内，当汽车的平均速度 \(v\) 为多少时，车流量最大？最大车流量为多少？（精确到 \(0.1\text{ 千辆/小时}\)）
\item 若要求在该时段内车流量超过 \(10\text{ 千辆/小时}\)，则汽车的平均速度应在什么范围内？
\end{enumerate}
\end{problem}

\begin{solution}
\begin{enumerate}
\item 因为 \(v>0\)，有
\[
v+\frac{1600}{v}\geq2\sqrt{1600}=80
\]
于是
\[
v^2+3v+1600=v\left(v+3+\frac{1600}{v}\right)\geq83v
\]
分母为正，所以
\[
y=\frac{920v}{v^2+3v+1600}\leq\frac{920}{83}
\]
基本不等式取等号当且仅当 \(v=\dfrac{1600}{v}\). 结合 \(v>0\)，得 \(v=40\). 因此平均速度为 \(40\text{ 千米/小时}\) 时车流量最大，最大车流量为
\[
\frac{920}{83}\approx11.1\text{ 千辆/小时}
\]
\item 因为分母为正，车流量超过 \(10\text{ 千辆/小时}\) 等价于
\[
\frac{920v}{v^2+3v+1600}>10
\]
即
\[
920v>10v^2+30v+16000
\]
整理并除以正数 \(10\)，得
\[
v^2-89v+1600<0
\]
由于
\[
v^2-89v+1600=(v-25)(v-64)
\]
抛物线开口向上，故不等式的解为 \(25<v<64\). 这已经满足题设 \(v>0\)，所以所求速度范围为 \((25,64)\text{ 千米/小时}\).
\end{enumerate}
\end{solution}

\begin{problem}
现有一组互不相同且从小到大排列的数据：\(a_0\)，\(a_1\)，\(a_2\)，\(a_3\)，\(a_4\)，\(a_5\)，其中 \(a_0=0\). 为提取反映数据间差异程度的某种指标，今对其进行如下加工：记 \(T=a_0+a_1+\cdots+a_5\)，\(x_n=\dfrac{n}{5}\)，\(y_n=\dfrac{1}{T}(a_0+a_1+\cdots+a_n)\)，作函数 \(y=f(x)\)，使其图象为逐点依次连接点 \(P_n(x_n,y_n)\)（\(n=0\)，\(1\)，\(2\)，\(\cdots\)，\(5\)）的折线.

\begin{enumerate}
\item 求 \(f(0)\) 和 \(f(1)\) 的值；
\item 设 \(P_{n-1}P_n\) 的斜率为 \(k_n\)（\(n=1\)，\(2\)，\(3\)，\(4\)，\(5\)），判断 \(k_1\)，\(k_2\)，\(k_3\)，\(k_4\)，\(k_5\) 的大小关系；
\item 证明：\(f(x_n)<x_n\)（\(n=1\)，\(2\)，\(3\)，\(4\)）.
\end{enumerate}
\end{problem}

\begin{solution}
由题设 \(0=a_0<a_1<a_2<a_3<a_4<a_5\)，所以各项均为正，且 \(T>0\). 记
\[
R_n=a_0+a_1+\cdots+a_n
\]
则 \(y_n=\dfrac{R_n}{T}\)，并且折线在节点处满足 \(f(x_n)=y_n\).

\begin{enumerate}
\item 因为 \(x_0=0\)，且 \(R_0=a_0=0\)，所以 \(P_0=(0,0)\)，从而 \(f(0)=0\). 又因为 \(x_5=1\)，且 \(R_5=T\)，所以 \(P_5=(1,1)\)，从而 \(f(1)=1\).
\item 相邻两点的横坐标之差为 \(x_n-x_{n-1}=\dfrac15\)，纵坐标之差为 \(y_n-y_{n-1}=\dfrac{R_n-R_{n-1}}{T}=\dfrac{a_n}{T}\). 因此
\[
k_n=\frac{y_n-y_{n-1}}{x_n-x_{n-1}}=\frac{\frac{a_n}{T}}{\frac15}=\frac{5a_n}{T}
\]
因为 \(T>0\)，且 \(a_1<a_2<a_3<a_4<a_5\)，所以
\[
k_1<k_2<k_3<k_4<k_5
\]
\item 对 \(n=1\)，\(2\)，\(3\)，\(4\)，由 \(a_{n+1}\)，\(\ldots\)，\(a_5\) 中每一项都大于 \(a_n\)，有
\[
T-R_n=a_{n+1}+\cdots+a_5>(5-n)a_n
\]
另一方面，\(a_0=0\)，且 \(a_1\)，\(\ldots\)，\(a_n\) 均不大于 \(a_n\)，所以 \(R_n\leq na_n\). 因此
\[
nT-5R_n=n(T-R_n)-(5-n)R_n>n(5-n)a_n-(5-n)R_n\geq0
\]
所以 \(nT>5R_n\). 再由 \(f(x_n)=y_n=\dfrac{R_n}{T}\) 及 \(x_n=\dfrac n5\)，得到
\[
f(x_n)=\frac{R_n}{T}<\frac n5=x_n
\]
\end{enumerate}
\end{solution}
